\section{Pendapatan Asli Daerah (PAD)}
\label{sec:pad}

PAD merupakan akumulasi dari pos penerimaan pajak yang terdiri atas pajak daerah dan retribusi daerah, pos penerimaan non pajak berupa penerimaan hasil perusahaan milik daerah, serta pos penerimaan investasi serta pengelolaan sumber daya alam \cite{nasir2019pad}. PAD bertujuan memberikan kewenangan kepada Pemerintah Daerah untuk mendanai pelaksanaan otonomi daerah sesuai dengan potensi daerah sebagai wujud desentralisasi. Semakin besar dana PAD yang diperoleh oleh daerah akan sebanding dengan laju pembangunan di daerah tersebut. 

Sesuai dengan Undang-Undang No. 33 tahun 2004 tentang perimbangan keuangan antara pusat dan daerah pasal 6 bahwa sumber Pendapatan Asli Daerah meliputi : pajak daerah, retribusi daerah, hasil pengelolaan kekayaan daerah yang dipisahkan, dan lain-lain PAD yang sah. Pendapatan daerah lainnya yang sah yang disebutkan sebelumnya berasal dari hasil penjualan kekayaan daerah yang tidak dipisahkan, jasa giro, pendapatan bunga, keuntungan selisih nilai tukar rupiah terhadap mata uang asing, dan komisi, potongan, ataupun bentuk lain sebagai akibat dari penjualan dan/atau pengadaan barang dan/atau jasa oleh Daerah \cite{uu33_2004}. 

Pajak yang dipungut oleh pemerintah kabupaten/kota berdasarkan Undang-undang (UU) Nomor 1 Tahun 2022 tentang Hubungan Keuangan antara Pemerintah Pusat dan Pemerintah Daerah terdiri atas:

\begin{enumerate}
    \item PBB-P2 (Pajak Bumi dan Bangunan Perdesaan dan Perkotaan)\\
    Contoh: Pajak tahunan yang dibayarkan atas kepemilikan rumah, apartemen, atau ruko di kota tersebut.
    \item BPHTB (Bea Perolehan Hak atas Tanah dan Bangunan)\\
    Contoh: Biaya pajak yang dibayarkan satu kali pada saat membeli rumah baru atau bekas dan melakukan proses balik nama sertifikat.
    \item PBJT (Pajak Barang dan Jasa Tertentu)\\
    Ini adalah gabungan dari beberapa pajak, contoh paling umum adalah:
    \begin{itemize}
        \item Pajak Restoran: Pajak 1\% yang tercantum di struk saat makan di restoran atau kafe.
        \item Pajak Hotel: Pajak yang dibebankan pada tagihan saat menginap di hotel.
        \item Pajak Parkir: Pajak yang pengelola parkir bayarkan (dan biasanya dibebankan ke pengunjung) saat parkir di mal atau gedung perkantoran.
        \item Pajak Hiburan: Pajak pada tiket saat menonton bioskop atau konser.
    \end{itemize}
    \item Pajak Reklame\\
    Contoh: Pajak yang dibayar oleh sebuah toko atau perusahaan untuk memasang papan nama toko, billboard, atau neon box di pinggir jalan.
    \item PAT (Pajak Air Tanah)\\
    Contoh: Pajak yang dibayar oleh pabrik, hotel, atau mal yang menggunakan pompa air sumur bor (air tanah) untuk kegiatan operasionalnya, bukan menggunakan air PDAM.
    \item Pajak MBLB (Pajak Mineral Bukan Logam dan Batuan)\\
    Contoh: Pajak yang dibayar oleh perusahaan tambang pasir, kerikil, atau batu kapur saat mereka mengambil material tersebut dari alam.
    \item Pajak Sarang Burung Walet\\
    Contoh: Pajak yang dibayar oleh seorang pengusaha yang memanen dan menjual sarang burung walet dari gedung atau gua.
    \item Opsen PKB (Opsen Pajak Kendaraan Bermotor)\\
    Contoh: Tambahan biaya (pungutan ekstra) yang nanti akan dikenakan di atas pajak STNK tahunan yang dibayarkan untuk mobil atau motor.
    \item Opsen BBNKB (Bea Balik Nama Kendaraan Bermotor) \cite{uu1_2022}.
    Contoh: Tambahan biaya (pungutan ekstra) yang akan dikenakan di atas biaya normal saat membeli kendaraan baru atau melakukan balik nama kendaraan bekas.
\end{enumerate}